\documentclass[10pt,twocolumn,letterpaper]{article}

\usepackage{geometry}
\geometry{letterpaper, margin=1in}

\usepackage{times}
\usepackage{microtype}
\usepackage{booktabs}
\usepackage{tabularx}
\usepackage{hyperref}
\usepackage{amsmath}
\usepackage{listings}
\usepackage{xcolor}

\urlstyle{same}
\hypersetup{
    colorlinks=true,
    linkcolor=blue,
    filecolor=magenta,      
    urlcolor=blue,
    pdftitle={MinEnglish Language Specification},
    pdfauthor={Dan Micsa, PhD},
}

\lstset{
    basicstyle=\ttfamily\small,
    breaklines=true,
    frame=single,
    backgroundcolor=\color{gray!10},
    keywordstyle=\bfseries,
    commentstyle=\itshape
}

\title{\Huge\textbf{MinEnglish Language Specification (v2.0)}\\\LARGE A Zero-Redundancy Phonetic Communication Protocol}
\author{\textbf{Dan Micsa, PhD}\\MIT License (Appendix)}
\date{March 2026}

\begin{document}

\maketitle

\begin{abstract}
The MinEnglish language is a formal communication protocol engineered to maximize Shannon information density while minimizing cognitive load and parser ambiguity. By systematically eliminating morphological irregularities and orthographic redundancies inherent in natural languages, MinEnglish achieves high computational consistency. This specification outlines the phonological, morphological, and syntactic architecture of version 2.0.
\end{abstract}

\section{Design Principles and Underlying Rationale}

To optimize human-to-human and human-to-machine transmission, MinEnglish operates on five absolute constraints:

\begin{itemize}
    \item \textbf{Phonetic Isomorphism:} 1-to-1 grapheme-phoneme correspondence eliminates the cognitive overhead of orthographic memorization and provides predictable tokenization for machine parsers.
    \item \textbf{Morphological Parsimony:} Pluralization, verb conjugation, and declension are replaced by explicit numerical and temporal prefixes, removing ambiguity in morphological boundaries.
    \item \textbf{ASCII Strictness:} Operates exclusively within the standard US keyboard namespace without diacritics. The letters \textit{q}, \textit{w}, and \textit{y} are omitted to balance acoustic distinctness with key space economy.
    \item \textbf{Operator Integration:} Mathematical symbols (\verb|*|, \verb|~|, \verb|>|, \verb|<|, \verb|:|) function as first-class syntactic tokens, compressing complex semantic modifiers into single-byte representations.
    \item \textbf{Absolute Regularity:} A zero-exception framework guarantees that programmatic rules apply universally, ensuring high predictability for human acquisition and natural language processing models.
\end{itemize}

\section{Phonology and Orthography}

\subsection{Vowels}
To ensure acoustic discernibility across diverse linguistic backgrounds, the vowel space is constrained to five cardinal points: \verb|a| (\textit{cat}), \verb|e| (\textit{bed}), \verb|i| (\textit{bit}), \verb|o| (\textit{hot}), and \verb|u| (\textit{but}).

\subsection{Long Vowels and Diphthongs}
Diphthongs and elongated vowels are encoded as contiguous grapheme pairs, preserving the 1-to-1 spelling mapping without requiring macron diacritics: \verb|ei| (\textit{day}), \verb|ii| (\textit{see}), \verb|ai| (\textit{like}), \verb|ou| (\textit{go}), \verb|uu| (\textit{food}), \verb|au| (\textit{how}), \verb|oi| (\textit{boy}), \verb|ar| (\textit{car}), \verb|or| (\textit{for}), \verb|er| (\textit{her}).

\subsection{Consonants}
Consonantal mappings prioritize acoustic clarity while resolving historical orthographic collisions.

\begin{table}[h!]
\centering
\begin{tabularx}{\linewidth}{lXl}
\toprule
\textbf{Graph.} & \textbf{Application Rule} & \textbf{Example} \\
\midrule
\textbf{c} & Soft, unaspirated velar plosive & cat, bac \\
\textbf{k} & Hard, aspirated velar plosive & king, kic \\
\textbf{ci} & Voiceless palato-alveolar affricate & ciald \\
\textbf{u} & Semivowel substitution for \textit{w} & uoter \\
\textbf{i} & Semivowel substitution for \textit{y} & ies \\
\textbf{t} & Alveolar plosive / Voiceless dental fricative & tinc \\
\textbf{d} & Alveolar plosive / Voiced dental fricative & dat \\
\textbf{sh} & Standard /$\int$/ sound & ship \\
\textbf{zh} & Standard /$\dots$/ sound (vision) & vizhun \\
\bottomrule
\end{tabularx}
\caption{Non-standard Consonant Mappings}
\end{table}

\textbf{Aspiration Differentiation (\texttt{k} vs \texttt{c}):} The aspirated /k\textsuperscript{h}/ is strictly represented by \texttt{k}, whereas the unaspirated variant uses \texttt{c}. Following sibilants (e.g., /s/), where aspiration naturally diminishes, \texttt{c} is mandated (e.g., \textit{scan}).

\subsection{Orthographic Normalization Constraints}
Historical orthographic artifacts must be algorithmically stripped. Silent vowels/consonants are nullified. The \texttt{-tion} suffix is phoneticized to \texttt{-shun}. Soft \texttt{c} regressively substitutes to \texttt{s}.

\subsection{Lexical Prosody Marker}
Because MinEnglish eliminates redundant double consonants and silent terminal vowels, lexical stress can become ambiguous. The apostrophe (\texttt{'}) is instituted as an optional prosody marker, affixed immediately preceding the stressed vowel: \texttt{ecsp'ensiv}, \texttt{compi'uuter}.

\section{Morphology and Syntax}

\textbf{Core Constraint: Morphological Invariance.} Lexical items never undergo morphological derivation or inflection. Grammatical relationships are managed exclusively through standardized operators and syntactic ordering.

\subsection{Nouns and Quantitative Prefixes}
MinEnglish mandates explicit numerical prefixing. Morphological plurals are abolished. 
\begin{itemize}
    \item \texttt{1}: Unitary (\texttt{1cat})
    \item \texttt{3}: Exact (\texttt{3cat})
    \item \texttt{*}: Universal (\texttt{*cat})
    \item \texttt{\~{}5}: Approximate (\texttt{\~{}5cat})
    \item \texttt{0}: Null set (\texttt{0cat})
\end{itemize}

\subsection{The Literalization Operator (\texttt{\_})}
To arrest lexicalization (where compound nouns fuse away from their root meanings over time), the Literalization Operator (\texttt{\_}) is introduced. Hyphens (\texttt{-}) denote conceptual, culturally unified compounds (\texttt{haus-fuud} = restaurant). Underscores (\texttt{\_}) force the computation of strict compositional roots (\texttt{haus\_fuud} = a house literally constructed from food).

\subsection{Proper Noun Encapsulation}
Proper Nouns (specific entities or brands) retain their canonical native orthography to preserve search retrieval. First-letter capitalization acts as an escape character, suspending local phonetic rules for the token duration (e.g., \texttt{1person Isabella}, \texttt{1cuntri France}).

\subsection{Fractional and Decimal Operators}
Spoken mathematics is standardized via the \texttt{/} (slais) and \texttt{.} (point) operators.
\begin{itemize}
    \item \textbf{Fractions:} \texttt{1/2} (uan slais tuu) = One-half. 
    \item \textbf{Decimals:} \texttt{0.5} (ziro point faiv).
\end{itemize}

\subsection{Verbs --- Temporal Mechanics}
Time processing targets an absolute mathematical timeline via temporal prefixing. A verb absent a modifier defaults to the continuous present.

\begin{table}[h!]
\centering
\begin{tabularx}{\linewidth}{lX}
\toprule
\textbf{Operator} & \textbf{Temporal Value} \\
\midrule
\textit{(null)} & Present continuous (\texttt{iit}) \\
\texttt{1} & Present discrete (\texttt{i1iit}) \\
\texttt{*} & Habitual continuity (\texttt{*iit}) \\
\texttt{-0s} & Immediate past vector (\texttt{-0s sii}) \\
\texttt{-1d} & Past vector, one day (\texttt{-1d iit}) \\
\texttt{:5Y} & Duration scalar, 5 years (\texttt{:5Y studi}) \\
\texttt{+1d} & Future vector, one day (\texttt{+1d gou}) \\
\texttt{2026-12} & Absolute timestamp \\
\bottomrule
\end{tabularx}
\caption{Temporal and Vector Prefixing}
\end{table}

\textbf{Temporal Stacking:} MinEnglish resolves deep time by permitting left-to-right stacking of temporal operators, enabling precise algebraic timeline construction (e.g., \texttt{i-0s -1d iit} = "I had already eaten yesterday").

\subsection{Modifiers and Scalar Intensity}
Adjectives post-modify nouns. Adverbs post-modify verbs.
Intensity is modified via combinators:
\begin{itemize}
    \item \texttt{>big}: Positive intensity (very big)
    \item \texttt{>>big}: Extreme positive intensity
    \item \texttt{<hot}: Negative intensity (slightly hot)
\end{itemize}

\subsection{Pronominal Indexing}
Pronouns operate as highly compressed, single-character spatial indices: \texttt{i} (1st sing.), \texttt{u} (2nd sing.), \texttt{h} (3rd masc.), \texttt{s} (3rd fem.), \texttt{t} (inanimate). Plurality logic inherits canonical numerical prefixing (e.g., \texttt{5i} = we five, \texttt{*t} = inanimate collective).

\textbf{Syntactic Agglutination:} The pronominal index agglutinates physically to the succeeding temporal operator, eliminating whitespace parsing overhead (\texttt{h1laic}, \texttt{u:2h uorc}).

\subsection{Prepositional Object Marking}
Like Romance languages, the default object requires no marker. If the syntactic sequence implies the noun could be a secondary action, the dative marker \texttt{tu} forces the parser into an objective state (\texttt{i1laic tu 1cat}).

\subsection{Clause Anchors (\texttt{co} / \texttt{oc})}
When nesting two or more clauses, speakers must use \texttt{co} (clause open) to start the subset and \texttt{oc} (clause close) to return to the parent set, mitigating short-term memory (phonological loop) bounds.

\subsection{Echo-Tagging (Noisy Protocol)}
For critical transmissions in high-loss environments, Shannon redundancy may be artificially inflated to 100\% by "echoing" the quantitative prefix as a full word after the root constraint (e.g., \texttt{3boi trii 1uoc nau}).

\section{Formal Syntax (Backus-Naur Form)}
The absolute regularity of MinEnglish allows it to be algorithmically validated.

\begin{lstlisting}[language=SQL]
<sentence> ::= [<subject>] <verb_block> [<obj>]
<subject>  ::= <noun_phrase> | <pronoun_phrase>
<obj>      ::= <noun_phrase> | <pronoun_phrase>

<noun_phrase> ::= <qty> <root> [<adjective>]
<pronoun_phrase> ::= [<qty>] <root>
<adjective> ::= [<intensity>] <adj_root>

<verb_block> ::= [<temporal>] <root> [<adverb>]
<temporal> ::= <time> | <duration> | <unix>

<clause> ::= "co" <connect> <sentence> "oc"
\end{lstlisting}

\section{Academic Data Lexicon}

\begin{table}[h!]
\centering
\begin{tabularx}{\linewidth}{lX}
\toprule
\textbf{Standard English} & \textbf{MinEnglish Lexicon} \\
\midrule
Database & \texttt{daatabeis} \\
Variable & \texttt{ting-ceinj} \\
Hypothesis & \texttt{aidia-test} \\
Algorithm & \texttt{rul-mat} \\
Equation & \texttt{balans-mat} \\
System & \texttt{net-ting} \\
Energy & \texttt{pauer-mun} \\
Cell & \texttt{bocs-laif} \\
Evolution & \texttt{ceinj-taim} \\
\bottomrule
\end{tabularx}
\caption{Computational Compounds}
\end{table}

\section{Comparative Information Density}

Transmission efficiency against a 35-sentence standard English corpus yields an internal \textbf{Average Information Density Gain of 35.0\%} over natural American English. Maximum localized compression exceeds 63.4\% when isolating absolute numerical timestamping syntax against conversational equivalents.

\section{Limitations: Biological Friction}
The primary identified limitation of MinEnglish exists externally to its syntactic structure; the limitation is biological. Mammalian communicative pathways evolved to facilitate social cohesion utilizing ambiguity to reduce friction. MinEnglish operates on an absolute algebraic paradigm. Widespread neurological adoption would prompt measurable psychological resistance, as the protocol disables the linguistic mechanisms humans rely on to process and mitigate sociological complexity.

\end{document}
